\chapter{KMS States and Thermal Correlators}
\label{ch:kms-states-thermal-correlators}

\ConventionsRealTime

Thermal equilibrium is a property of a state together with a time evolution.
The KMS condition expresses equilibrium without assuming a finite-volume
Hamiltonian trace formula, and therefore applies directly to
infinite-volume thermal QFT.

\section{The KMS Condition}

\begin{definition}[\(C^\ast\)-dynamical system]
\label{def:cstar-dynamical-system}
A \(C^\ast\)-dynamical system is a pair \((\Obs,\alpha_t)\), where
\(\Obs\) is a \(C^\ast\)-algebra and \(\alpha_t\) is a strongly continuous
one-parameter group of \(\ast\)-automorphisms of \(\Obs\).
\end{definition}

\begin{definition}[KMS state]
\label{def:thermal-volume-kms-state}
Let \(\beta>0\).  A state \(\omega:\Obs\to\C\) is a \(\beta\)-KMS state for
\(\alpha_t\) if, for every pair \(A,B\) in a norm-dense \(\alpha_t\)-analytic
subalgebra, there is a function \(F_{A,B}(z)\), holomorphic on
\[
  0<\operatorname{Im}z<\beta,
\]
continuous on the closed strip, such that
\[
  F_{A,B}(t)=\omega(A\alpha_t(B)),
  \qquad
  F_{A,B}(t+\ii\beta)=\omega(\alpha_t(B)A)
\]
for all real \(t\).
\end{definition}

If \(\Obs\) is represented on a finite-dimensional Hilbert space and
\(\alpha_t(A)=\ee^{\ii Ht}A\ee^{-\ii Ht}\), then
\[
  \omega_\beta(A)
  =
  \frac{\operatorname{Tr}(\ee^{-\beta H}A)}
       {\operatorname{Tr}(\ee^{-\beta H})}
\]
satisfies the KMS condition by cyclicity of the trace.  Conversely, the KMS
condition remains meaningful in infinite-volume regimes where the Gibbs trace
formula has no finite partition function.

\section{Spectral Form}

Assume the GNS representation of a KMS state has a self-adjoint generator
implementing \(\alpha_t\) on a common domain for two operators \(A,B\).  Define
the thermal Wightman functions
\[
  G_{AB}^{>}(t)=\omega(A(t)B(0)),
  \qquad
  G_{AB}^{<}(t)=\omega(B(0)A(t)),
\]
where \(A(t)=\alpha_t(A)\).  The KMS boundary condition gives
\[
  G_{AB}^{>}(t+\ii\beta)=G_{AB}^{<}(t).
\]
For tempered boundary values, Fourier transformation yields the detailed
balance relation
\[
  \widetilde G_{AB}^{<}(\omega)
  =
  \ee^{-\beta\omega}\widetilde G_{AB}^{>}(\omega).
\]
The spectral density
\[
  \rho_{AB}(\omega)
  =
  \widetilde G_{AB}^{>}(\omega)
  -
  \widetilde G_{AB}^{<}(\omega)
\]
therefore determines both thermal two-point functions:
\[
  \widetilde G_{AB}^{>}(\omega)
  =
  \frac{\rho_{AB}(\omega)}{1-\ee^{-\beta\omega}},
  \qquad
  \widetilde G_{AB}^{<}(\omega)
  =
  \frac{\ee^{-\beta\omega}\rho_{AB}(\omega)}
       {1-\ee^{-\beta\omega}} .
\]
For fermionic operators the graded KMS condition inserts the fermion-parity
sign in the boundary equation and gives the corresponding Fermi factors.

\section{Linear Response And Retarded Correlators}

Let \(B(t)\) be an observable coupled to a compactly supported real source
\(f(t)\) by the time-dependent generator
\[
  H_f(t)=H-f(t)B(t)
\]
in a finite system, or by the corresponding derivation on an analytic
subalgebra in the algebraic setting.  To first order in \(f\), Duhamel's
formula gives the change of the expectation of another observable \(A(t)\):
\[
  \delta\omega(A(t))
  =
  \int_{-\infty}^{t} ds\, f(s)\,
  G^R_{AB}(t-s),
  \qquad
  G^R_{AB}(t)
  =
  -\ii\,\theta(t)\,\omega([A(t),B(0)]) .
\]
Indeed, expanding the interaction-picture evolution operator gives
\[
  U_f(t,-\infty)
  =
  1+\ii\int_{-\infty}^{t}ds\,f(s)B(s)+O(f^2),
\]
and inserting \(U_f^{-1}A(t)U_f\) in the unperturbed KMS state produces the
commutator displayed above.  The retarded distribution is causal by
construction: its support in the time variable lies in \(t\ge0\).  Its
Fourier transform is the boundary value of a function holomorphic in the
upper half \(\omega\)-plane whenever the time-growth hypotheses needed for
the transform hold.

For a translation-invariant thermal QFT, the retarded two-point function of
local densities is
\[
  G^R_{AB}(\omega,\mathbf k)
  =
  -\ii\int_0^\infty dt\int d^{D-1}x\,
  \ee^{\ii\omega t-\ii\mathbf k\cdot\mathbf x}
  \omega([A(t,\mathbf x),B(0,\mathbf0)]) .
\]
The spectral function associated to \(G^R\) is
\[
  \rho_{AB}(\omega,\mathbf k)
  =
  -2\,\operatorname{Im}G^R_{AB}(\omega+\ii0,\mathbf k)
\]
with this sign convention.  The KMS detailed-balance relation above then
expresses the symmetrized fluctuation spectrum in terms of the dissipative
part of the retarded response.  This is the precise QFT input behind
transport formulae.

\section{Conserved Densities And Hydrodynamic Variables}

Hydrodynamics begins with conserved local currents in the microscopic state
space.  Let \(T^{\mu\nu}\) be the stress tensor and let
\(J_A^\mu\) be conserved currents for ordinary internal symmetries:
\[
  \partial_\mu T^{\mu\nu}=0,
  \qquad
  \partial_\mu J_A^\mu=0 .
\]
In a homogeneous isotropic thermal state with unit future timelike vector
\(u^\mu\), temperature \(T\), and chemical potentials \(\mu_A\), the one-point
functions define thermodynamic densities by
\[
  \omega(T^{\mu\nu})
  =
  \varepsilon u^\mu u^\nu+p\Delta^{\mu\nu},
  \qquad
  \omega(J_A^\mu)=n_A u^\mu,
  \qquad
  \Delta^{\mu\nu}=\eta^{\mu\nu}+u^\mu u^\nu .
\]
The variables \(\varepsilon,n_A,u^\mu\) become hydrodynamic fields only after
one assumes a family of states whose long-wavelength correlators are
controlled by slowly varying local values of these densities.  The derivative
expansion is then a statement about constitutive relations:
\[
  T^{\mu\nu}
  =
  T^{\mu\nu}_{(0)}+T^{\mu\nu}_{(1)}+T^{\mu\nu}_{(2)}+\cdots,
  \qquad
  J_A^\mu
  =
  J_{A,(0)}^\mu+J_{A,(1)}^\mu+J_{A,(2)}^\mu+\cdots,
\]
where the subscript records the number of derivatives acting on
\(T,\mu_A,u^\mu\) and background sources.

At first order, parity-even relativistic hydrodynamics in flat spacetime
contains shear viscosity \(\eta_{\rm sh}\), bulk viscosity \(\zeta\), and
charge-conductivity data.  In the shear channel, the Kubo formula follows
from the linear-response identity by coupling a metric perturbation
\(h_{xy}(t)\) to \(T^{xy}\).  With the retarded-sign convention above,
\[
  \eta_{\rm sh}
  =
  \lim_{\omega\downarrow0}
  \frac{\rho_{T^{xy}T^{xy}}(\omega,\mathbf0)}{2\omega}
  =
  -\lim_{\omega\downarrow0}
  \frac{\operatorname{Im}G^R_{T^{xy}T^{xy}}(\omega,\mathbf0)}{\omega}.
\]
This formula is a definition of the transport coefficient in terms of QFT
correlators once the zero-frequency limit exists.  Later chapters in this
volume must derive the constitutive relations, entropy-current constraints,
Schwinger--Keldysh effective action, anomaly-induced transport, and kinetic
limits from these correlator-level objects.

\begin{openproblem}[Thermal real-time reconstruction]
\label{op:thermal-real-time-reconstruction}
Give constructive hypotheses under which Euclidean finite-temperature
Schwinger functions on \(S^1_\beta\times\R^{D-1}\) reconstruct a real-time
KMS QFT with local observables, spectral functions, and transport
coefficients.  The required analytic continuation is more delicate than the
vacuum OS reconstruction because thermal correlators have periodic Euclidean
time and no vacuum spectrum condition.
\end{openproblem}
